\documentclass[a4paper,12pt]{article}
\usepackage[utf8]{inputenc}
\usepackage[T1]{fontenc}
\usepackage[ngerman]{babel}
\usepackage{geometry}
\geometry{margin=2.5cm}

\begin{document}

\section*{APEBench: A Benchmark for Autoregressive Neural Emulators of PDEs}
\textbf{Autoren:} Felix Koehler, Simon Niedermayr, Rüdiger Westermann, Nils Thuerey \\
\textbf{Konferenz:} NeurIPS 2024 (Datasets and Benchmarks Track) \\
\textbf{Code:} \texttt{https://github.com/tum-pbs/apebench}

\subsection*{Zusammenfassung}

APEBench ist ein umfassendes Benchmark-Framework zur systematischen Evaluierung autoregressiver neuronaler Emulatoren für zeitabhängige partielle Differentialgleichungen (PDEs). Das Framework basiert auf JAX und integriert einen hocheffizienten, differenzierbaren Pseudo-Spektral-Löser auf Basis von Exponential Time Differencing Runge-Kutta (ETDRK) Methoden, der 46 verschiedene PDE-Dynamiken in ein, zwei und drei räumlichen Dimensionen abdeckt.

\paragraph{Lösungsansatz und Solver-Integration.}
Ein zentrales Alleinstellungsmerkmal von APEBench gegenüber bestehenden Benchmarks wie PDEBench oder PDEArena ist die enge Integration des numerischen Referenzlösers direkt in den Trainings- und Evaluierungsprozess. Die unterstützten semi-linearen PDEs besitzen die Form
\[
\partial_t u = \mathcal{L}u + \mathcal{N}(u),
\]
wobei der lineare Anteil $\mathcal{L}$ im Fourier-Raum diagonalisiert und damit exakt über eine Matrix-Exponentialfunktion integriert werden kann, während der nichtlineare Anteil $\mathcal{N}(u)$ durch Runge-Kutta-Stufen approximiert wird. Dies ermöglicht für lineare PDEs eine nahezu fehlerfreie Datengenerierung, was die Qualität der Trainingsgrundlage erheblich verbessert. Für nichtlineare Terme wird ein Pseudo-Spektral-Ansatz mit De-Aliasing nach Orszag (2/3-Regel) verwendet.

\paragraph{Schwierigkeitsidentifikatoren.}
Ein weiterer wichtiger Beitrag ist die Einführung standardisierter Schwierigkeitskennzahlen $\gamma_j$ für jede PDE-Dynamik, die direkt an klassische numerische Stabilitätskriterien geknüpft sind. Für Ableitungen $s$-ter Ordnung mit Koeffizient $a_s$ werden diese definiert als
\[
\gamma_s = \alpha_s N^s 2^{s-1} D, \quad \text{mit} \quad \alpha_s = \frac{a_s \Delta t}{L^s},
\]
wobei $N$ die Auflösung, $D$ die Raumdimension und $L$ der Domänenumfang sind. Für $s=1$ entspricht $\gamma_1$ genau der CFL-Zahl der Advektionsgleichung. Durch diese Parametrisierung lässt sich eine PDE-Konfiguration eindeutig identifizieren und ihre Emulierschwierigkeit intuitiv einschätzen.

\paragraph{Architekturvergleich.}
APEBench evaluiert fünf repräsentative neuronale Architekturen: einfache Faltungsnetzwerke (Conv), Residualnetzwerke (ResNet), U-Nets, Dilated ResNets und den Fourier Neural Operator (FNO). Zentrale Erkenntnisse sind:
\begin{itemize}
  \item \textbf{ResNets} erweisen sich über nahezu alle PDE-Typen und räumliche Dimensionen hinweg als konsistent stärkste Architektur, insbesondere durch ihre Skip-Connections, die mit steigender Dimensionalität immer wichtiger werden.
  \item \textbf{Lokale Faltungsarchitekturen} (Conv, ResNet) sind auf Probleme mit begrenztem Einflussbereich gut geeignet, scheitern jedoch, wenn die effektive Rezeptivfeldgröße die physikalisch notwendige unterschreitet.
  \item \textbf{FNO} profitiert von seiner globalen Rezeptivfeldstruktur und zeigt besonders bei Navier-Stokes-Strömungen (Kolmogorov-Flow) starke Ergebnisse. Bei Reaktions-Diffusions-Gleichungen hingegen leidet er unter einem Low-Frequency-Bias, da er hochfrequente Muster nur indirekt über Nichtlinearitäten repräsentieren kann.
  \item \textbf{U-Net} zeigt in höheren Dimensionen (3D) gute Leistung dank seines globalen Rezeptivfeldes auf grober Skala.
\end{itemize}

\paragraph{Trainingsstrategien und Unrolling.}
APEBench legt besonderen Wert auf die systematische Untersuchung des sogenannten \textit{Unrolled Training}, bei dem das Netzwerk während des Trainings mehrere Zeitschritte autoregressiv vorwärts rolliert und der Fehler über die gesamte Rollout-Länge akkumuliert. Die Autoren formalisieren verschiedene Trainingsparadigmen durch eine einheitliche Taxonomie:
\begin{itemize}
  \item \textbf{One-Step Supervised Training} ($T = B = 1$): Klassisches Ein-Schritt-Training, das keine differenzierbare Simulation erfordert.
  \item \textbf{Supervised Unrolling} ($T = B$): Alle Referenztrajektorien werden vorab berechnet.
  \item \textbf{Diverted-Chain Unrolling} ($T > B = 1$): Der Referenzlöser wird bei jedem Schritt auf die aktuelle Netzwerkausgabe angewandt und muss daher differenzierbar sein.
\end{itemize}
Die Experimente belegen, dass längeres Unrolling die Langzeitstabilität der Emulatoren systematisch verbessert, allerdings auf Kosten der Kurzzeitgenauigkeit. Die Diverted-Chain-Methode kombiniert dabei die Vorteile beider Extrema: Sie erreicht hohe Langzeitgenauigkeit ohne signifikante Einbußen in der Kurzzeitleistung.

\paragraph{Neuronale Hybrid-Emulatoren.}
Schließlich unterstützt APEBench auch Korrektur-Tasks, bei denen ein grobes numerisches Verfahren $\tilde{\mathcal{P}}_h$ als Prädiktor fungiert und ein neuronales Netz die verbleibenden Fehler korrigiert. Diese Konfiguration entspricht konzeptuell dem Residuallern-Ansatz und zeigt, dass insbesondere ResNets von einer solchen Arbeitsteilung profitieren: Wenn der grobe Löser bereits 50\,\% der Schwierigkeit abdeckt ($\tilde{\gamma}_1 = 0.5\,\gamma_1$), reduziert sich der Rollout-Fehler erheblich.

\subsection*{Bezug zur Masterarbeit}

Die Ergebnisse von APEBench sind für die vorliegende Masterarbeit in mehrfacher Hinsicht relevant, auch wenn sich die Problemstellung grundlegend unterscheidet: Während APEBench zeitabhängige PDE-Dynamiken emuliert, befasst sich die Masterarbeit mit der stationären Vorhersage von Stokes-Strömungsfeldern um sedimentierende Kristalle in magmatischen Systemen.

\paragraph{Architekturwahl: U-Net vs.\ FNO.}
Die in APEBench beobachteten Stärken und Schwächen der untersuchten Architekturen sind direkt auf den Anwendungsfall der Masterarbeit übertragbar. Der Befund, dass U-Nets in höherdimensionalen Szenarien durch ihr hierarchisches Rezeptivfeld gut abschneiden, unterstützt deren Einsatz für die Vorhersage räumlich ausgedehnter Geschwindigkeitsfelder. Gleichzeitig bestätigt APEBench die Stärke des FNO für globale Strömungsmuster -- eine Eigenschaft, die für die langreichweitigen Wechselwirkungen in Mehrpartikel-Stokes-Strömungen relevant ist. Dass ResNets durchweg kompetitiv bleiben und besonders dateneffizient sind (Konvergenz bereits mit wenigen Trainingsfällen), ist mit Blick auf die begrenzte Datenbasis von LaMEM-Simulationen in der Masterarbeit wichtig.

\paragraph{Residuallern-Ansatz und Hybrid-Emulatoren.}
Der in APEBench untersuchte Korrektur-Ansatz -- bei dem ein neuronales Netz Abweichungen von einem groben Hintergrundfeld korrigiert -- entspricht konzeptuell dem in der Masterarbeit implementierten Residuallernverfahren, bei dem analytische Stokes-Einzelpartikellösungen als Prädiktor dienen und das Netzwerk nur die Mehrpartikel-Wechselwirkungskorrekturen $\Delta v_\text{learned}$ erlernen muss. Die APEBench-Ergebnisse bestätigen, dass eine solche Hybridstrategie den Lernaufwand signifikant reduziert.

\paragraph{Physikalische Konsistenz und Divergenzfreiheit.}
APEBench betont explizit die Analogie zwischen neuronalen Emulatoren und klassischen numerischen Verfahren. Der Befund, dass Emulat\-oren bei zunehmend längerem Unrolling numerisch stabileren Differenzschemata ähneln, spiegelt die in der Masterarbeit verfolgte Stream-Function-Strategie wider: Durch die explizite Parameterisierung über eine Stromfunktion $\psi$ ist die Bedingung $\nabla \cdot \mathbf{v} = 0$ per Konstruktion exakt erfüllt -- eine im Sinne der APEBench-Ergebnisse elegante Möglichkeit, physikalische Konsistenz ohne Soft-Constraint-Verlustterme zu garantieren.

\paragraph{Generalisierung und Skalierung.}
Das zentrale Forschungsziel der Masterarbeit -- die Generalisierung von einem auf $N$ Kristalle trainierten Modell auf $N+m$ Kristalle -- findet im APEBench-Rahmen eine Analogie in der Generalisierung über unterschiedliche PDE-Schwierigkeiten $\gamma$. APEBench zeigt, dass Architekturen mit unzureichendem Rezeptivfeld bei steigender Schwierigkeit versagen, was die Bedeutung geeigneter Normierungsstrategien und der Wahl einer physikalisch motivierten Eingabedarstellung für die Generalisierung in der Masterarbeit unterstreicht.

\end{document}